\documentclass[11pt,a4paper]{article}  % Déclare la classe du document.

\newcommand{\chaptermark}[2]{\og #2 \fg{} (#1)} % ligne uniquement pour éviter une erreur en cas d'article
\include{../1ere-preambule}

\usepackage{epic}   % extension et simplification de  l'environnement picture
\usepackage{graphicx}    % permet d'inclure des graphique externe dans un document
\usepackage{arydshln}   %permet d'avoir des lignes en pointillés

\newcommand{\lignes}[1]{
	\setlength{\unitlength}{1mm}
	\vskip 0,4cm\noindent
	\multido{\i=0+1}{#1}{%
		\noindent
		\begin{picture}(150,5)
			\linethickness{0,005mm}
			\put(-5,0){\dashline{1}(175,5)(5,5)}
		\end{picture}\vskip 0,001cm
	}
	\vskip 0cm
}

\rhead{DS n$^{\circ}02$ $-$ Automatismes $-$ 05/11/2025}

\newcommand{\va}[1]{\lvert \ #1 \ \rvert}
\newcommand{\vaf}[1]{\left\lvert \ #1\  \right\rvert}

\begin{document}
	\graphicspath{{images/}}
	
	\noindent Nom: ${\scriptstyle  .........................................................}$ \hskip 1.8cm Prénom: ${\scriptstyle  ...................................................}$ 
	
	\vspace{.5cm}
	
	\begin{center}
		
		
		\fbox{
			\begin{minipage}[]{15cm}
				\begin{center}
					\LARGE{\rule{0.5cm}{0pt}\rule[-0.25cm]{0pt}{0.9cm}
						05/11/2025 $-$ DS n$^{\circ}02$ $-$ Automatismes\\
						Proportion et pourcentage 
					}
				\end{center}
			\end{minipage}
		}
	\end{center}

\begin{center}\textbf{55min - Barème indicatif}\end{center}
\begin{center}\textbf{Les élèves avec un tier-temps ne traitent pas les questions avec  le symbole \ding{46}}\end{center}
\begin{center}\textbf{Calculatrice non autorisée.\\Les résultats devront être justifiés (à l'aide de calcul ou de propriétés).}\end{center}



\exerciceDS{(Calculer un pourcentage) $-$ 6 points}

Calculer:
\begin{multicols}{2}
	\begin{enumMet}
		\item $25\ \%$ de $300$ ;
		\item \ding{46} $30\ \%$ de $220$;
		\item $0,5\ \%$ de $64\ 400$;
		\item $140\ \%$ de $20$.
		\item $5\ \%$ de $380$.
		\item $13\ \%$ de $500$.
	\end{enumMet}
\end{multicols}

\exerciceDS{(Calculer une proportion ou un effectif) $-$ 5 points}
\begin{enumMet}
	\item Calculer $p$ lorsque $n_S = 300$ et $n_E = 1\ 500$.\\
	Écrire le résultat sous la forme d'un pourcentage.
	\item \ding{46} Calculer $p$ lorsque $n_S = 30$ et $n_E= 9\ 000$.\\
	Écrire le résultat sous la forme d'un pourcentage.
	\item Calculer $n_S$ lorsque $p = 0,12$ et $n_E= 5\ 000$.
	\item Calculer $n_E$ lorsque $p = 0,30$ et $n_S = 180$.
	\item Sachant que $30 \%$ d'une somme $S$ vaut $4\ 500$ \EUR{}, calculer $S$.
\end{enumMet}

\exerciceDS{ $-$ Taux d'évolution entre $Q_{1}$ et $Q_{2}$ $-$ 4 points}
Dans chacun des cas suivants , calculer l'un des trois nombres $Q_{1}, Q_{2}$ ou $t$, connaissant les deux autres. De plus donnner le taux d'évolution sous forme de pourcentage et indiquer s'il s'agit d'une baisse ou d'une hausse.
	\begin{multicols}{2}
		\begin{enumMet}
		\item \ding{46} $Q_{1}=4 ; Q_{2}=2$.
		\item $Q_{1}=10 ; Q_{2}=14$.
		\item $Q_{1}=30 ; t=-0,30$.
		\item $Q_{2}=260 ; t=0,3$.
	\end{enumMet}
	\end{multicols}

\exerciceDS{$-$ 2 points}

Compléter les phrases suivantes.
\begin{multicols}{2}
	\begin{enumMet}
	\item Augmenter de $35 \%$, c'est multiplier par \phantom{$\dfrac{1}{1}$}......
	\item Diminuer de $25 \%$, c'est multiplier par \phantom{$\dfrac{1}{1}$}......
	\item Multiplier par 6 c'est augmenter de \phantom{$\dfrac{1}{1}$}.......\%.
	\item Multiplier par 0,55 c'est diminuer de \phantom{$\dfrac{1}{1}$}.......\%.	
\end{enumMet}
\end{multicols}

\exerciceDS{ $-$ 4,5 points}

\begin{enumMet}
	\item Une personne paie, pour un groupe, une note de restaurant qui s'élève à 360 \EUR{}, avec le service compris de $20 \%$. Quel est le prix des repas sans le service?
	\item \ding{46} Un commerçant calcule ses prix de vente en prenant un bénéfice de $10 \%$ sur ses prix d'achat. Quel est le prix d'achat d'un article qu'il a vendu $330$\EUR{}.
	\item Le prix d'un article soldé est de $180$ \EUR{}. L'étiquette indique \guillemets{$-40 \%$}. Calculer le prix de l'article avant les soldes.
\end{enumMet}


\end{document}

